\lyseis{Γ2.2}{Εύρος, κορυφές και απόμακρες τιμές}
\sol{1} α) $12-3=9$. β) $6-(-4)=10$. γ) $15-15=0$. δ) $2{,}5-0{,}5=2$.
\sol{2} α) Το $3$ (τρεις φορές). β) Το $4$ και το $6$ (από δύο φορές), δύο κορυφές. γ) Το «μπλε» (τρεις φορές).
\sol{3} α) Λάθος: μεγαλύτερη μείον μικρότερη τιμή, άρα δεν είναι αρνητικό. β) Σωστό: μεγαλύτερη και μικρότερη τιμή ταυτίζονται. γ) Λάθος: μπορεί να υπάρχουν πολλαπλές κορυφές. δ) Σωστό. ε) Σωστό.
\sol{4} α) Το $48$. β) Το $3$. γ) Καμία: όλες οι τιμές είναι κοντά η μία στην άλλη.
\sol{5} α) $2+7+3+6+1=19$ μαθητές. β) $50-0=50$ λεπτά. γ) Ναι, δύο: οι κλάσεις $10$--$20$ και $30$--$40$. δ) $2+7=9$ μαθητές.
\sol{6} Άθροισμα $-2+3+0+5+2+8-2=14$, μέση τιμή $14:7=2$ $^\circ$C. Εύρος $8-(-2)=10$ $^\circ$C. Κορυφή το $-2$ (δύο φορές).
\sol{7} Α: μέση τιμή $7$, εύρος $9-5=4$. Β: μέση τιμή $7$, εύρος $13-1=12$. Μεγαλύτερη μεταβλητότητα έχει το Β, αφού έχει το ίδιο κέντρο αλλά μεγαλύτερο εύρος.
\sol{8} Αγγελική: $\dfrac{85}{5}=17$ δευτ., εύρος $18-16=2$. Πολυξένη: $\dfrac{85}{5}=17$ δευτ., εύρος $20-14=6$. Ίδιος μέσος όρος, αλλά πιο σταθερή είναι η Αγγελική (μικρότερο εύρος). Τον καλύτερο χρόνο έκανε η Πολυξένη ($14$ δευτ., όσο πιο μικρός ο χρόνος τόσο καλύτερα).
\sol{9} α) Με επτά μέρες: άθροισμα $245$, μέση τιμή $35$. Στη σειρά: $5,36,38,40,40,42,44$, διάμεσος $40$. Εύρος $44-5=39$. Χωρίς τη βροχερή: άθροισμα $240$, μέση τιμή $40$. Στη σειρά: $36,38,40,40,42,44$, διάμεσος $40$. Εύρος $44-36=8$. β) Η διάμεσος: έμεινε $40$. Η μέση τιμή άλλαξε κατά $5$ και το εύρος από $8$ έγινε $39$.
\sol{10} Εύρος $165-140=25$ cm. Κέντρα: $142{,}5$, $147{,}5$, $152{,}5$, $157{,}5$, $162{,}5$. Γινόμενα: $285$, $737{,}5$, $1220$, $630$, $162{,}5$. Άθροισμα $3035$. Μέση τιμή $3035:20=151{,}75$ cm.
\sol{11} Δύο κορυφές: ο βαθμός $14$ και ο βαθμός $15$ (από $6$ μαθητές ο καθένας). Εύρος $17-12=5$. Άθροισμα βαθμών $24+39+84+90+48+34=319$, μέση τιμή $319:22=14{,}5$.
\sol{12} Πίνακας: τιμή $1$: συχνότητα $1$, $2$: $2$, $3$: $2$, $4$: $1$, $5$: $4$, $6$: $2$ (σύνολο $12$). Κορυφή το $5$. Εύρος $6-1=5$. Άθροισμα $47$, μέση τιμή $47:12\approx3{,}9$.
\sol{13} Παράδειγμα: $1,5,7,7,11$. Εύρος $11-1=10$, διάμεσος $7$ (η μεσαία), κορυφή $7$ (εμφανίζεται δύο φορές). Υπάρχουν και άλλες απαντήσεις.
\sol{14} α) Νέα δεδομένα $5, 8, 11$: μέση τιμή $8$ (αυξήθηκε κατά $3$), εύρος $6$ (ίδιο). β) Νέα δεδομένα $4, 10, 16$: μέση τιμή $10$ (διπλασιάστηκε), εύρος $12$ (διπλασιάστηκε). Παρατήρηση: όταν προσθέτουμε σταθερό αριθμό, η μέση τιμή μετακινείται, αλλά το εύρος μένει ίδιο (τα δεδομένα μετατοπίζονται χωρίς να απλώνονται περισσότερο). Όταν πολλαπλασιάζουμε, αλλάζουν και τα δύο.
\sol{15} Νέα μέση τιμή $\dfrac{30+40}{7}=10$. Νέα διάμεσος: στη σειρά $3,4,5,5,6,7,40$, η μεσαία είναι $5$. Η απόμακρη τιμή διπλασίασε τη μέση τιμή και δεν άλλαξε τη διάμεσο.
\sol{16} α) Δύο: οι κλάσεις $10$--$20$ και $40$--$50$. β) Στο πανηγύρι ήρθαν κυρίως παιδιά και έφηβοι και γονείς (ενήλικες): δύο ομάδες ηλικιών. γ) Κέντρα $5,15,25,35,45$. Άθροισμα $5\cdot2+15\cdot7+25\cdot1+35\cdot3+45\cdot6=10+105+25+105+270=515$. Πλήθος $19$. Μέση τιμή $515:19\approx27{,}1$ έτη. Δεν περιγράφει σχεδόν κανέναν: στην κλάση $20$--$30$ ήρθε μόνο ένας.
\sol{17} Ελεύθερη εργασία. Π.χ. $18,20,17,19,21$ λεπτά: μέση τιμή $19$, διάμεσος $19$, εύρος $21-17=4$.
\sol{18} Ελεύθερη εργασία. Η κορυφή είναι το χρώμα που εμφανίζεται συχνότερα. Μέση τιμή δεν έχει νόημα, γιατί τα χρώματα δεν είναι αριθμοί και δεν μπορούμε να τα προσθέσουμε.
